\section{Thesis Structure}
In essence, the thesis is aiming to solve one fundamental question:

\textit{“What is the sensitivity of the climate response to the superposition of two or more aircraft exhaust plumes, under a variety of atmospheric conditions?”} 

As such, the document is structured as follows: Chapter 2 introduces the science behind aviation climate impact, and explores a wide range of literature pertaining to aircraft emissions, their plume-scale effects, and the spatio-temporal sensitivity of the atmospheric response. Chapter 3 investigates the feasibility of aircraft exhaust plume intersections in modern airspace, through consideration of real-world air traffic data and rudimentary dispersion calculations, to estimate degree of plume overlap in high-density regions. In chapter 4, the potential to offset climate change through formation flying is investigated in more detail, with a focus on gas-phase chemical kinetics associated with \ce{NO_x} emissions and their varying impact on ozone formation. This leads to the development of a localised plume chemistry modelling framework called the Gridded Plume Analysis Tool (GPAT) in chapter 5. In this chapter, the methodology behind a novel aircraft plume analysis tool is outlined and supporting models and methods are introduced. GPAT is also validated against its predecessor for mass conservation and consistency across implementations of the photochemical trajectory model. In chapter 6, this modelling framework is then utilised in a formation flight study, to investigate the most critical parameters affecting the severity of the \ce{NO_x} saturation phenomenon, and the benefits and drawbacks of formation flying in the context of aviation climate impact mitigation.